\documentclass{article}

\usepackage[utf8]{inputenc}     
\usepackage[T1]{fontenc}
\usepackage[magyar]{babel}

\usepackage{amsmath}
\usepackage{amsthm}
\usepackage{amssymb}

\title{Moduláris Aritmetika}
\author{Simon László}
\date{Legutolsó frissítés: 2026. Szeptember 15.}

\newtheorem{theorem}{Tétel}
\newtheorem{definition}{Definíció}[section]

\begin{document}

\maketitle
\tableofcontents
\newpage

\section{Relációk}
Egy reláció lehet szimmetrikus, reflexív, és tranzitív.\footnote{Egy olyan relációt, amelyre mindhárom telyesül, ekvivalenciarelációnak nevezzük}
\begin{definition}[Szimmetrikus]
    Egy relációt (legyen $\sim$) szimmetrikusnak nevezzük, ha minden $a,b$-re teljesül, hogy $a \sim b \iff b \sim a$.
\end{definition}
\begin{definition}[Reflexív]
    Egy relációt (legyen $\sim$) reflexívnek nevezzük, ha minden $a$-ra igaz, hogy $a \sim a$.
\end{definition}
\begin{definition}[Tranzitív]
    Egy relációt (legyen $\sim$) tranzitívnek nevezzük, ha minden $a, b, c$-re igaz, hogy $(a \sim b) \vee (b \sim c) \implies (a \sim c)$.
\end{definition}

\begin{center}
\begin{tabular}{c|c|c|c|}
    R. & Sz. & Tr. & Reláció \\
    \hline
    + & + & + & = \\
    \hline
    + & + & - & $(a \sim a) \vee (a \sim (a-1)) \vee (a \sim (a+1))$ \\
    \hline
    + & - & + & Oszthatóság \\
    \hline
    + & - & - & $(a \sim a)\vee (a \sim (a+1))$ \\
    \hline
    - & + & + & nincs! \\
    \hline
    - & + & - & $\perp$ \\
    \hline 
    - & - & + & < \\
    \hline
    - & - & - & $a \sim (a+1)$ \\
\end{tabular}
\end{center}

\section{Kongruenciák}
\begin{definition}[Kongruencia]
    $a \equiv b \pmod{m} \iff m \mid b - a, m > 0$
\end{definition}

\begin{theorem}[A kongruencia ekvivalenciareláció]
\end{theorem}
\begin{proof} $ $\\
\begin{enumerate}
    \item Reflexivitás \\
    $\forall a$ $a \equiv a \pmod{m} \iff m \mid a - a = 0$\\
    $\implies$ reflexív.
    \item Szimmetria \\
    $\forall a,b$ $a \equiv b \pmod{m} \implies b \equiv a \mod{m} \iff m \mid b - a \iff m \mid a - b$ \\
    $\implies$ szimmetrikus.
    \item Tranzitivitás \\
    $\forall a,b,c$ $a \equiv b \pmod{m} \vee b \equiv c \pmod{m} \iff m \mid b - a \vee m \mid c - b \implies m \mid c - a$ \\
    $\implies$ tranzitív.
\end{enumerate}
\end{proof}

Mivel a kongruencia egy ekvivalenciareláció, az alaphalmazból ekvivalencia-osztályokat hoz létre. \\
Egy ekvivalencia-osztályra a következőek teljesülnek: \\
\begin{enumerate}
    \item Az alaphalmaz részhalmaza
    \item Disjunktok
    \item Uniójuk az alaphalmaz
\end{enumerate}


\end{document}